\section{WC 27/1/25}

\subsection*{Overview:}
This week I took to Figma to implement my research into some design concepts I can begin to build within my webapp

\subsection*{Achievements, Thoughts and Ideas:}
With a new understanding of the principles behind a both intuitive and enjoyable Ui experience, I set work on recreating this for my own project. 

Firstly, I took to paper (well digital paper) to draw some basic ideas. Whilst the professional tool I previously researched (Figma), is very powerful at creating polished designs, it cant beat the freethinking and fast changing benefits of a blank canvas. I thought about the different goals a user would have when accessing my project and laid my views out to help them achieve these goals in the most time and thought efficient way. My ideas and brainstorming can be found here: \nameref{fig:brainstormUI}.

With these ideas In mind I then created a new Figma project and added a variety of 'frames' which each represent the view a user might have. Following this I began adding a number of rectangles to represent various aspects of my software. It was at this point I remembered back to my Figma research and about the magical tool that was components libraries, using a library of premade components would enable me to 'drag and drop' relevant elements into my design vastly improving the speed of my prototyping.

I soon found the Official Figma library created by the team at Shadcdn ( the UI library I have chosen to build my software), excellent I thought as these components would perfectly match what I can easily create and make possible as I code my prototype. Unfortunately this wasn't quite the case. Scrolling through the list of elements included in the library a number seemed to be missing that I had previously identified and had drawn my eye as I was choosing Shadcdn (including the sidebar components which I had read much praise about from developers on platforms such as StackOverflow and Reddit). It seems this much just be an internal delay from the team as the missing parts are relatively new however I couldn't afford to wait for this crucial update as the design of my UI was now constraining other aspects of my project.

Abandoning the official route I decided to search through community made alternatives. After looking through a couple which also had challenges such as a diversion from the official colour palettes / design standards and more outdated components I found one which perfectly fit the bill (\href{https://www.figma.com/design/JgunNpiqpurqV1ANxtD466/shadcn-ui-component-library--Community---Copy-?node-id=112-11267}{Link}) This library had every aspect I needed, and I excitedly hit the glowing add to project button… until I faced another setback.

Since my previous research of Figma (and the time since the tutorial I watched was made and I viewed it) the software had changed pricing model to focus more towards teams and enterprises. This new paywall locked the 'plugins' feature (which wrapped around libraries) behind a paywall. This was disheartening however with the understanding I had of the power of Figma as a platform I decided to take a look at the pricing models to see if I could fund my own account for a short time, fortunately I discovered they had a very generous student discount of 100\% to entice aspiring designers into their product suite prior to employment. After submitting proof of education status I received my shiny new account the next day and set to work.

To the right -> you can see the results of my first few iterations of my app dashboard page. I've focused primarily on the layout of core features such as settings, note finding etc

Within the EPQ lesson this week I spent some time reviewing my peers work and having my own work reviewed. This was an exciting opportunity to gain feedback and critique and whilst I am very proud of the marks I was given, I've also seen some areas I need to further improve (mostly regarding the documentation of my sources and timelines as they adapt over time)

\subsection*{Power of Three:}
\textbf{This week's Tasks laid out per day}

\begin{tabularx}{\textwidth}{
  >{\raggedright\arraybackslash}X  % First column expands
  >{\centering\arraybackslash}p{1.5cm}  % Middle column has fixed width
  >{\raggedright\arraybackslash}X  % Third column expands
}
\toprule
\textbf{Goal} & \textbf{Achieved?} & \textbf{Notes} \\
\midrule


\rowcolor{gray!10}
\multicolumn{3}{l}{\textbf{Monday}} \\
Create new Figma project & \ding{51} & \\
Register for student premium account & \ding{51} & \\
Organise component libraries ready for design layout & \ding{51} &  \\



\rowcolor{gray!10}
\multicolumn{3}{l}{\textbf{Tuesday}} \\
Build Wireframes of page layouts & \ding{51} & \\
Select components to fill layouts& \ding{51} & \\
Quick A/B testing with peers on design challenges & \ding{51} &  \\



\rowcolor{gray!10}
\multicolumn{3}{l}{\textbf{Thursday}} \\
Complete layouts  & \ding{51} & \\
Select 'break points' for reactive content to screen size & \ding{51} & \\
Polish sections with basic prototyping & \ding{51} & \\  


\rowcolor{gray!10}
\multicolumn{3}{l}{\textbf{Thursday}} \\
Review mark schemes for EPQ coursework& \ding{51} & \\
'mark' EPQ partners current work  & \ding{51} & \\
Reflect on own mark awarded by partner & \ding{51} &  Please find here: \nameref{Evaluation:SpringPeerReview}\\

\bottomrule
\end{tabularx}


\subsection*{Plans for next week:}  
Next week I hope to take these Ui designs and implement them into another prototype which I can show to peers for feedback
